\subsection{کالبدشکافی ماژول \lr{salt}}

\subsubsection*{۱. فلسفه طراحی، پارادایم معماری و تمثیل ملموس}
معماری \lr{salt} با هدف هم‌راستاسازی طیفی وزن‌های کم‌رتبه با جهت‌های اصلی واریانس اطلاعاتی مدل مبنا توسعه یافته است. این ماژول پایه‌های تبدیل را روی بردارهای ویژه متناظر با بزرگ‌ترین مقادیر تکین قفل می‌کند تا همگرایی تنظیم دقیق با کمترین پارامتر ممکن محقق گردد.

\textbf{تمثیل ملموس و مثال کاربردی:}
کوک کردن سیم‌های یک تار یا پیانو بر مبنای تشدید آکوستیک چوب کاسه ساز: به جای سفت کردن تصادفی سیم‌ها، کشش سیم‌ها دقیقاً طوری تنظیم می‌شود که بسامدهای ارتعاشی آن‌ها روی هارمونیک‌های رزونانس طبیعی کاسه صوتی قفل شوند تا صدایی رسا و پرقدرت با کمترین مصرف انرژی صوتی طنین‌انداز گردد.

\subsubsection*{۲. فرمول‌بندی و استنتاج عمیق ریاضی (Mathematical Derivation)}
ماتریس‌های نگاشت کم‌رتبه $\mathbf{A} \in \mathbb{R}^{r \times D}$ و $\mathbf{B} \in \mathbb{R}^{D \times r}$ بر پایه تجزیه مقادیر تکین ماتریس پیش‌آموزش‌دیده $\mathbf{W}_0 = \mathbf{U} \boldsymbol{\Sigma} \mathbf{V}^\top$ راه‌اندازی می‌شوند:
\begin{equation}
\mathbf{A}_0 = \mathbf{U}_r^\top, \quad \mathbf{B}_0 = \boldsymbol{\Sigma}_r \mathbf{V}_r^\top
\end{equation}
فرآیند به‌روزرسانی گرادیان به زیرفضای پوشش داده‌شده توسط مقادیر تکین ماکزیمم مقید می‌شود:
\begin{equation}
\mathbf{W}_{\text{eff}} = \mathbf{W}_0 + \frac{\alpha}{r} (\mathbf{X} \mathbf{B}) \mathbf{A}
\end{equation}
این هم‌راستاسازی اولیه مانع از آشفتگی در فاز تنظیم دقیق و افت دقت مدل می‌گردد.

\subsubsection*{۳. کالبدشکافی الگوریتمی و جریان داده (Algorithmic Tensor Trace)}
در پیاده‌سازی سورس پایتون در فایل ماژول \lr{salt.py}، پردازش در متد \lr{forward} به صورت مرحله‌به‌مرحله و گام‌به‌گام صورت می‌پذیرد:
\begin{enumerate}
\item تجزیه طیفی وزن لایه توجه و استخراج $r=16$ بردار ویژه برتر ماتریس مبنا.
\item مقداردهی اولیه شاخه‌های کم‌رتبه موازی بر اساس پایه‌های طیفی استخراج‌شده.
\item محاسبه پروجکشن ورودی به موازات لایه منجمدشده اصلی.
\item اعمال ضریب مقیاس‌بندی تطبیقی و ترکیب تانسورهای دو مسیر با جمع برداری.
\item ادغام تحلیلی وزن‌ها در زمان استنتاج جهت حذف کامل هرگونه سربار تأخیر.
\end{enumerate}

\subsubsection*{۴. تحلیل همجوشی سخت‌افزاری و ترافیک حافظه (Hardware \& Roofline Analysis)}
از دیدگاه مدل سقف کارایی (\lr{Roofline Model}) و معماری پردازنده‌های گرافیکی مدرن (مانند \lr{NVIDIA Ada / RTX 4070}):
\begin{itemize}
\item \textbf{صفر شدن سربار محاسباتی در فاز استنتاج:} با ادغام ماتریس‌های کم‌رتبه در وزن پایه پیش از استنتاج، هیچ سربار اضافی تحمیل نمی‌شود.
\item \textbf{کاهش ۹۹ درصدی پارامترهای گرادیان:} حجم حافظه مصرفی بهینه‌ساز در فاز آموزش به شدت افت می‌کند.
\end{itemize}

\subsubsection*{۵. ارزیابی تجربی، بنچمارک و توجیه علمی نتایج}
بر اساس آزمون‌های ثبت‌شده در خط لوله ارزیابی پایان‌نامه:
\begin{itemize}
\item \textbf{وضعیت رسمی آزمون کیفی:} \lr{PASS}
\item \textbf{دقت ثبت‌شده روی بنچ‌مارک:} 80.28 درصد
\item \textbf{شاخص ارزیابی آماری (\lr{F1 Score}):} 78.9 درصد
\item \textbf{اوج مصرف حافظه گرافیکی (\lr{Peak VRAM}):} 750.0 مگابایت (در برابر ۴۲۰۰ مگابایت مدل مبنا)
\item \textbf{نرخ توان عملیاتی (\lr{Throughput}):} 36,100 توکن بر ثانیه
\end{itemize}

\textbf{توجیه علمی و علت‌شناسی جذاب نتایج:}
کسب دقت 80.28 درصد و کنترل اوج مصرف رم در سطح 750.0 مگابایت (صرفه‌جویی 82.1 درصدی نسبت به ۴۲۰۰ مگابایت پایه) با سرعت 36,100 توکن/ثانیه، نشان می‌دهد که هم‌راستاسازی بردارهای تطبیق با طیف ویژه ماتریس پایه، از سرگردانی گرادیان در زیرفضاهای پوچ ممانعت می‌نماید.

\subsubsection*{۶. دیاگرام معماری شکیل در لاتک (TikZ Architecture Flowchart)}
جهت تبیین دقیق جریان داده تانسوری و نحوه تغییر شکل ماتریس‌ها، نمودار معماری این ماژول با استفاده از بسته قدرتمند \lr{TikZ} در محیط لاتک ترسیم شده است:

\begin{center}
\begin{tikzpicture}[scale=0.82, every node/.style={scale=0.82}, 
    node distance=1.1cm,
    auto,
    block/.style={rectangle, draw=blue!70, fill=blue!12, rounded corners=4pt, text centered, minimum height=0.85cm, font=\small},
    spec/.style={rectangle, draw=purple!70, fill=purple!15, rounded corners=4pt, text centered, minimum height=0.85cm, font=\small},
    outblock/.style={rectangle, draw=teal!80, fill=teal!15, rounded corners=4pt, text centered, minimum height=0.85cm, font=\small}
]
    \node [block] (input) {ورودی $\mathbf{X}$};
    \node [spec, right=1.0cm of input] (align) {هم‌راستاسازی طیفی};
    \node [spec, right=1.0cm of align] (lora) {آداپتور کم‌رتبه $\mathbf{B}\mathbf{A}$};
    \node [outblock, right=1.0cm of lora] (out) {وزن ادغام‌شده};

    \draw [->, thick, draw=blue!80] (input) -- (align);
    \draw [->, thick, draw=purple!80] (align) -- (lora);
    \draw [->, thick, draw=purple!80] (lora) -- (out);
\end{tikzpicture}
\end{center}

\subsubsection*{۷. پیوست تصویری و پرامپت تولید هوش مصنوعی}
برای ترسیم تجسم سه‌بعدی و گرافیک صنعتی این مکانیزم با ابزارهای هوش مصنوعی، پرامپت انگلیسی زیر با دقت تدوین شده است:

\vspace{0.2cm}
\begin{latin}
\footnotesize
\textbf{AI Technical Visualization Prompt:}\\
"A 3D scientific visualization of Spectral Alignment and Low-Rank Tuning (SALT). Harmonic spectral resonance curves locking into alignment with structural crystalline lattice bars of a neural network, glowing amber and cyan tones, 8k resolution --ar 16:9"
\end{latin}

\begin{figure}[H]
    \centering
    \framebox[0.9\textwidth]{
    \begin{minipage}{0.85\textwidth}
    \vspace{1.8cm}
    \centering
    \textbf{[کادر الحاق تصویر گرافیکی سه‌بعدی ماژول \lr{salt}]}\\[0.2cm]
    \small (بر اساس پرامپت فنی فوق تولید و در این کادر جایگزین می‌شود)
    \vspace{1.8cm}
    \end{minipage}
    }
    \caption{شماتیک و جریان داده تانسوری در ماژول \lr{salt}}
    \label{fig:salt}
\end{figure}

\vspace{0.5cm}
\hrule
\vspace{0.5cm}
